\documentclass[12pt]{article}
\usepackage[utf8]{inputenc}
\usepackage[margin=0.75in]{geometry}
\usepackage{amsmath}
\usepackage{amssymb}
\usepackage{amsthm}
\usepackage{graphicx} % Required for inserting images
\usepackage{array}
\usepackage{tikz-cd}


\title{MAT1100 Homework 1}
\author{Aaron Avram}
\date{September 19 2026}

\begin{document}
\newtheorem{definition}{Definition}
\newtheorem{theorem}{Theorem}
\newtheorem{claim}{Claim}
\newtheorem{corollary}{Corollary}

\maketitle

\begin{align*}
    \textbf{Question 4}
\end{align*}
\begin{enumerate}
    \item[(a)]
    \begin{proof}
        To show $H$ is not normal in $G$, consider $r^{-1} \in G$ and $rs \in H$
        then
        \[ r^{-1}(rs)r = sr = r^{-1}s \]
        which is not in $H$, hence $H$ is not invariant under conjugation meaning it is not normal.
        To show that 
        \[ (aH)(bH) = (ab)H \]
        is not a well defined operation, it suffices to show that given $a', b'$, other representatives of $aH$ and $bH$
        
        \[ (a'b')H \neq (ab)H \]
        to  do this take $a = b = s$ and $a' = b' = s \cdot rs$ which is also a representative of $sH$
        as $rs \in H$. Now note that
        \[ ab = s^2 = e \]
        and
        \[ a'b' = s \cdot rs \cdot s \cdot rs = sr^2s = r^{-2}s^2 = r^2 \]
        and clearly $r^2 \notin H$ therefore
        \[ r^2H \neq H \]
        which implies the above operation is not well defined.
    \end{proof}

    \item[(b)]
    \begin{proof}
        To show that $N$ is normal recall that by question $3$, $N = Z(D_8)$
        and therefore for all $x \in D_8$ and $g \in N$
        \[ xg = gx \implies xgx^{-1} = g \in N \]
        thus $N$ is normal.
        The cosets of $N$ in $G$ are:
        \begin{align*}
            N = \{ e, r^2 \} \\
            rN = \{ r, r^3 \} \\
            sN = \{ s, sr^2 \} \\
            srN = \{ sr, sr^3 \} \\
        \end{align*}
    \end{proof}

    \item[(c)]
    \begin{proof}
        If $D_8/N$ were cyclic, it would be generated by a single element
        of order $|D_8|/|N| = 4$. Thus to show $D_8/N$ is not cyclic, it suffices
        to show that it contains no element of order greater than 2. To show that
        a given element $aN \in D_8/N$ has order at most 2, it suffices to pick
        a representative $g$ of $aN$ and show that $g^2 \in N$. By part $(b)$ the cosets of $N$ in $D_8$
        are
        \[\{ N, rN, sN, srN \} \]
        and fix representatives
        \[ \{ e, r, s, sr \} \]
        then
        \begin{align*}
            e^2 &= e \in N \\
            (r)^2 &= r^2 \in N \\
            s^2 &= e \in N \\
            (sr)^2 &= (sr) \cdot (r^{-1}s) = s^2 = e \in N \\
        \end{align*}
        hence all elements of $D_8/N$ have order at most 2 and thus the group cannot be cyclic
        as it has order 4.

        Now to prove that $D_8/N$ is abelian since every element of $D_8/N$ squares to the identity, it suffices to show that for any group $G$ if
        $x^2 = e$ for all $x \in G$ then $G$ is abelian. To do this, note that $x^2 = e$
        is equivalent to $x = x^{-1}$ and therefore if $a, b \in G$
        \[ ab = a^{-1}b^{-1} = (ba)^{-1} = ba \]
        where $(ba)^{-1} = ba \in G$ from the fact that every element in $G$ squares to the identity.
        Hence $G$ is abelian and therefore $D_8/N$ is abelian.
    \end{proof}
\end{enumerate}

\end{document}